\documentclass[../../main.tex]{subfiles}
\begin{document}

\begin{center}
  {\fontsize{48pt}{58pt}\selectfont \textbf{1991}\normalsize\textbf{年度}}
\end{center}

\begin{table}[htb]
  \centering
  \caption{出題テーマと難易度}
  \begin{tabular}{|c|p{9cm}|c|}
    \hline
    \textbf{問題番号} & \textbf{テーマ} & \textbf{難易度} \\
    \hline
    第1問 & 整数問題（階乗の末尾の0） & \\
    \hline
    第2問 & 空間図形と接平面 & \\
    \hline
    第3問 & 回転体の体積 & \\
    \hline
    第4問 & 三次関数と逆関数の積分 & \\
    \hline
    第5問 & 確率（三次方程式の整数解） & \\
    \hline
  \end{tabular}
\end{table}

\begin{table}[htb]
  \centering
  \caption{解答}
  \renewcommand{\arraystretch}{1.4}
  \begin{tabular}{|c|c|p{9cm}|}
    \hline
    \textbf{問題番号} & \textbf{小問} & \textbf{答え} \\
    \hline
    第1問 & (1) & $f(10)=2$，$f(100)=24$ \\
    \cline{2-3}
    & (2) & $\dfrac{1}{4}$ \\
    \hline
    第2問 & - & $2-\sqrt{2}\le k\le2+\sqrt{2}$ \\
    \hline
    第3問 & - & $V(a)=\dfrac{\pi}{3}\left(4+\dfrac{1}{2+a}+\dfrac{1}{2-a}\right)$ \\
    \hline
    第4問 & (2) & 最小値 $\dfrac{1}{12}$（$a=-2,b=\dfrac{4}{3}$） \\
    \hline
    第5問 & - & $\dfrac{19}{108}$ \\
    \hline
  \end{tabular}
\end{table}

\paragraph{試験形態}
5問．試験時間は要確認．

\end{document}
